% \begin{savequote}
% {\tiny Ein Mann ging von Jerusalem nach Jericho hinab und wurde von
% Räubern überfallen. Sie plünderten ihn aus und schlugen ihn nieder; 
% dann gingen sie weg und ließen ihn 
% halbtot liegen. Zufällig kam ein Priester denselben Weg
% herab; er sah ihn und ging weiter. Auch ein Levit kam zu
% der  Stelle; er sah ihn und ging weiter. Dann kam ein Mann
% aus  Samarien, der auf der Reise war. Als er ihn sah, hatte
% er  Mitleid, ging zu ihm hin, goss Öl und Wein auf seine
% Wunden  und verband sie. 
% 
% Dann hob er ihn auf sein Reittier,
% brachte  ihn zu einer Herberge und sorgte für ihn. Am
% anderen Morgen  holte er zwei Denare hervor, gab sie dem
% Wirt und sagte:  Sorge für ihn, und wenn du mehr für ihn
% brauchst, werde  ich es dir bezahlen, wenn ich wiederkomme.
% Was meinst du:  Wer von diesen dreien hat sich als der
% Nächste dessen erwiesen,  der von den Räubern überfallen
% wurde?  
% 
% Der Gesetzeslehrer antwortete: Der, der
% barmherzig\index{Barmherzigkeit} an  ihm gehandelt hat. Da
% sagte Jesus zu ihm: Dann geh und handle genauso!
% 
% Lk 10,25-37}
% % \qauthor{{\tiny Lk 10,25-37}}
% \end{savequote}
% 
\def\ChapterCode{HKL}
\chapter
[\CO{[HKL]} Herz-Kreislaufstörungen]
{Herz-Kreislaufstörungen}
\ChapterIntro
{%
Herz-Kreislauferkrankungen sind die häufigste
Todesursache.\footnote{Statistik Austria, Pressemitteilung
9.691-133/10; Statistisches Bundesamt Deutschland,
Pressemitteilung Nr.344 vom 15.09.2009} Fast jeder zweite
Todesfall geht auf eine Erkrankung des
Herz-Kreislauf-Systems zurück.  Mit steigendem Lebensalter
erhöht sich die Wahrscheinlichkeit an einer solchen zu
erkranken. 
Eine wichtige Gruppe stellen die erworbenen, durch den
Lebensstil bedingten Erkrankungen dar. Zu den wesentlichen
Risikofaktoren zählen Übergewicht, mangelnde Bewegung,
Rauchen, Bluthochdruck, erhöhte Blutfettspiegel und ein
erhöhter Blutzuckerspiegel.

Akute Notfälle ereignen sich oft nicht unvorhersehbar, oft
leidet der Patient vorab an \E{chronischen Varianten} einer
Herz-Kreislauferkrankung bzw. einer verwandten
Grunderkrankung, welche zu einem Notfallgeschehen führen
kann.
}
{%
\begin{description}
  \item[Maintainer:] Sebastian Gabriel
  \item[Autoren:] Diverse
  \item[Reviewer:] Standard-Reviewprozess
  \item[Version:] {\DocVersion} ({\DocVersionInfo})
  \item[SHA1:] {\DocGitCommit}
\end{description}

{\TxTeilkapitelAASS}
}


\nocite{Harrisons18De,Harrisons18DeKap224,Harrisons18DeKap225,Harrisons18DeKap226}


\Experimental{
\Layout{57}{}{}{}{%
\begin{tabularx}{\linewidth}{XXXX}\FontTableBase
&
\TabHeading{Chronische Variante bzw. Grunderkrankung}
&
\TabHeading{Akuter Notfall}
&
\TabHeading{Anm.}
\\\addlinespace\midrule\addlinespace

&
KHK, Angina pectoris
&
Akutes Koronarsyndrom (instabile Angina pectoris,
Herzinfarkt) 
&

\\\addlinespace

&
Hypertension
&
Hypertensive Krise, Hypertensiver Notfall
&

\\\addlinespace
&
pAVK
&
Artrieller Gefäßverschluss
&

\\\addlinespace

&
Vorhofflimmern, Störungen des Reizleitungssystems, \ldots
&
Tachykarde Attacke
&
\\\addlinespace

&
Herzinsuffizienz
&
Dekompensierte Herzinsuffizienz, kard. Lungenödem
&


% \\\addlinespace
% Grad
% &
% 
% &
% 
% &

\\\addlinespace

\end{tabularx}
}
{Beispiele für den Zusammenhang von chronischen und akuten
Herz-Kreislauf-Erkrankungen}
{}

} % STATUS-EXPERIMENTAL





\label{sec:herz-kreislauf-stoerungen}
\label{topic:herz-kreislauf-stoerungen}


\clearpage
% \section{Herz-Kreislauf-Stillstand}
% 
% Der Herz-Kreislauf-Stillstand wird
% unter
% \FullRef{topic:kreislaufstillstand},
% \FullRef{topic:reanimation-eh}, und
% \prettyref{chp:BLS} und \prettyref{chp:BLS} besprochen.



\section{Herzinsuffizienz}
\label{topic:herzinsuffizienz}
\index{Herzinsuffizienz}
\index{Herzinsuffizienz!kompensierte}
\index{Herzinsuffizienz!dekompensierte}
\index{Rechtsherzinsuffizienz}
\index{Linksherzinsuffizienz}
\index{Globalinsuffizienz}
\index{Insuffizienz!Herz-}


\DefinitionInPara{Eine \T{Herzinsuffizienz} bezeichnet eine 
unzureichende Auswurfleistung des Herzens, unabhängig von deren Ursache.}
Es sind viele Ursachen
möglich, manche Ursachen treten \E{akut} auf, manche sind
eher \E{chronischer} Natur.



\Layout{0}{Einteilung}
{Die Einteilung der Herzinsuffizienz kann nach betroffenen
Herzteil oder Fähigkeit des Körpers zur Kompensation
erfolgen:
\begin{itemize}
  \item Einteilung nach betroffenen \EE{Herzteil}:
  \begin{itemize}
    \item \T{Rechtsherzinsuffizienz}: Unzureichende Leistung des \E{rechten}
    Herzens,
    \item \T{Linksherzinsuffizienz}:
    Unzureichende Leistung des \E{linken} Herzens, und
    \item \T{Globalinsuffizienz}: Unzureichende Leistung beider Herzteile.
  \end{itemize}
  \item Einteilung nach \EE{Kompensation}: Abhängig davon, wie
  gut der Körper mit der
  Einschränkung zurecht kommt, unterteilt man:
  \begin{itemize}
    \item \T{Kompensierte Herzinsuffizienz}: Der Körper hat noch
    genug Reserven, kann gegensteuern und kommt mit der
    eingeschränkten Leistung aus.
    \item \T[Dekompensierte
    Herzinsuffizienz]{\E{De}kompensierte Herzinsuffizienz}:
    Der Körper kommt aufgrund des Fortschreitens der
    Erkrankung oder aufgrund eines erhöhten Bedarfs nicht mehr
    mit der Leistung aus. Es treten deutliche
    Symptome auf.
  \end{itemize}
\end{itemize}


}{
\begin{compactitem}
  \item Herzteil:
  \begin{compactitem}
    \item \E{Rechts}herzinsuffizienz
    \item \E{Links}herzinsuffizienz
    \item \E{Global}insuffizienz
  \end{compactitem}
  \item Kompensation:
  \begin{compactitem}
    \item Kompensierte Herzinsuffizienz
    \item \Ca{De}kompensierte Herzinsuffizienz
  \end{compactitem}
\end{compactitem}
}{}{}{}




\Layout{0}{{\TxSymptome} (allgemein) und Einschätzung}
{%
Es ist wichtig, den Schweregrad einer Herzinsuffizienz zu
beurteilen. Sie muss nicht dramatisch sein, für viele
Patienten ist sie ein \E{chronischer Zustand} mit dem sie
gut leben können, da sich der Körper mit der Zeit an die
abnehmende Leistungsfähigkeit gewöhnt.

Problematisch wird es jedoch, wenn die
\E{Kompensationsmechanismen} (köpereigene, Medikamente,
{\dots}) versagen und der Patient plötzlich eine Symptomatik
entwickelt. Man spricht dann von der \EE{dekompensierten}
Herzinsuffizienz. Man muss die Symptome generell \E{anhand
ihres Verlaufes beurteilen}: Bekannte, seit längerer Zeit
bestehende Symptome sind weniger alarmierend als plötzlich
auftretende Symptome. Plötzlich auftretende Symptome sind
oft \Ca{Alarmzeichen}, typischerweise Symptome wie
\EE{Atemnot}, \E{plötzlich auftretende
Ödeme} oder \E{akute Schmerzen}. 
Bei einem plötzlichen Leistungsabfall des \E{linken} Herzens
kommt es zu einem \E{Rückstau} des Blutes \E{in den
Lungenkreislauf}. Aufgrund der Stauung kommt es zum Austritt
von Flüssigkeit in das Zwischengewebe und es entsteht ein
\EE{(kardiales) Lungenödem}. Ein typisches Zeichen dafür ist
das \EE{brodelnde
Atemgeräusch}\index{Atemgeräusch!brodelndes - bei
Herzinsuffizienz}. Fällt die Leistung des \E{rechten}
Herzens, so kann man Beinödeme und/oder gestaute Halsvenen
beobachten.
}{%
\begin{compactitem}
  \item Chronisch: oft Gewöhnungseffekte, Kompensation
  \item Dekompensation: Plötzlich neu auftretende
  Symptome, Beschwerden
  \item \Ca{Alarmzeichen}
  \begin{compactitem}
    \item \EE{Atemnot}
    \item  {Akute Schmerzen}
    \item Stauung:
    \begin{compactitem}
      \item \EE{Lungenödem} mit \E{\enquote{brodelndem
      Atemgeräusch}} (\E{li.} Herz)
      \item \EE{Beinödeme}, gestaute Halsvenen (\E{re.}
      Herz)
    \end{compactitem}
  \end{compactitem}
\end{compactitem}
}{}{}{}



\clearpage

\Layout{0}{\TxSymptome: \E{Links}herzinsuffizienz}
{\label{topic:linksherzinsuffizienz}%
\index{Herzinsuffizienz!Links-}%
\index{Insuffizienz!Linksherz-}%
\index{Linksherzinsuffizienz}%
\index{Lungenödem!kardiales - bei Herzinsuffizienz}%
%
Die Linksherzinsuffizienz entsteht aufgrund mangelnder
Auswurfleistung des linken Herzens. Durch den \E{Rückstau
von Blut in den Lungenkreislauf} kann es zu einem
\EE{kardialen Lungenödem} kommen.
Der Patient klagt über Atemnot (\EE{Dyspnoe}). Oft berichten
die Patienten, dass sie nur mehr im Sitzen Luft bekommen,
sie können nicht flach liegen
(\Z{Orthopnoe}\ZFootnote{\E{Orthopnoe}: Der Patient muss
aufrecht sitzen um die Atemnot zu lindern.}).

Die Atemfrequenz und Herzfrequenz sind oft erhöht
(\EE{Tachypnoe},
\EE{Tachykardie}\index{Hyperventilation!bei
Herzinsuffizienz}), oft besteht ein \EE{Hustenreiz}.
Mitunter wirken die Patienten unruhig und ängstlich.
Grundsätzlich ist der Blutdruck vermindert (\EE{Hypotonie}).
Allerdings kann ein erhöhter Blutdruck bei einer
Blutdruckkrise (\EE{Hypertonie}, \E{hypertensive Krise}
(\prettyref{topic:hypertensive-krise}), \E{hypertensiver
Notfall} (\prettyref{topic:hypertensiver-notfall})) der
eigentliche \E{Auslöser} einer Herzinsuffizenz sein. Der
Blutdruck
kann dann natürlich erhöht sein. 

Ein besonders kritisches Symptom ist das \EE{Lungenödem}.
Wenn es sehr stark ausgeprägt ist, kann man es mit freiem
Ohr hören (\E{\enquote{brodelndes Atemgeräusch}}). Diese
Patienten sind als besonders kritisch einzuschätzen, sie
sind \E{lebensbedrohlich} krank. Das Herz verbraucht oft die
allerletzten Reserven, schon die Aufregung durch den
Transport im Stiegenhaus zum Auto kann das Herz endgültig in
den Herz-Kreislaufstillstand bringen. Es kommt immer wieder
vor, dass solche Patienten im Stiegenhaus oder im Auto
reanimationspflichtig werden.
Das kardiale Lungenödem wird unter
\prettyref{topic:lungenoedem} 
besprochen.

}{
\begin{compactitem}
  \item \EE{Dyspnoe}, Patient bekommt oft nur mehr im
    Sitzen Luft \Z{({Orthopnoe})}
  \item \EE{Hustenreiz}
  \item \EE{Tachypnoe}\index{Hyperventilation!bei Herzinsuffizienz}
  \item \EE{Tachykardie}
  \item \EE{Hypotonie} oder \EE{Hypertonie} (\Abk{RR}\,\Sign{↑}\Sign{↑} kann Auslöser
  sein)
  \item \EE{Lungenstauung/Lungenödem} (Blut staut sich
  zurück in den Lungenkreislauf)
  \begin{compactitem}
  \item Hörbares \Ca{Kardiales Lungenödem = lebensgefährlich
  krank!}
  \item Herz verbraucht die allerletzten Reserven!
\end{compactitem}
  \item Unruhe
\end{compactitem}
}{}{}{}

\Layout{0}{\TxSymptome: \E{Rechts}herzinsuffizienz}
{%
\label{topic:rechtsherzinsuffizienz}\index{Rechtsherzinsuffizienz}\index{Herzinsuffizienz!Rechts-}
\index{Insuffizienz!Rechtsherz-}%
Die Rechtsherzinsuffizienz entsteht aufgrund mangelnder
Auswurfleistung des rechten Herzens. Dementsprechend kommt
es zu einem Rückstau des Blutes in den Körperkreislauf und
damit zu Ödemen (Beine, \ldots) und gestauten Halsvenen.
Die Patienten müssen häufig nachts urinieren. \Z{Über längere
Zeit kann sich die Leber vergrößern und ein sog.
\enquote{Wasserbauch} (\Z{Aszites}) bilden, dabei wird durch
den Rückstau Wasser in den freien Bauchraum abgepresst.}
}{
\begin{itemize}
  \item \EE{Bein-}/\EE{Hautödeme}
  \item \EE{gestaute Halsvenen}
  \item Nächtliches Wasserlassen
\end{itemize}
}{}{}{}


\Layout{0}
{\TxABCDE}
{~
\begin{ABCDEText}
  \item[\ABCDE{Szeneüberblick}]
  \item[\ABCDE{Eindruck}] 
  \item[\ABCDE{Bewusstsein}]
  \item[\ABCDE{Hauptbeschwerde}] 
  \item[\ABCDE{A}] 
  \item[\ABCDE{B}] 
  \item[\ABCDE{C}] 
  \item[\ABCDE{D}] 
  \item[\ABCDE{E}] 
  \item[\ABCDE{=}] 
\end{ABCDEText}
}
{
\begin{ABCDESlide}
  \item[\ABCDE{Sz}] 
  \item[\ABCDE{Ei}] 
  \item[\ABCDE{Be}] 
  \item[\ABCDE{Hb}] 
  \item[\ABCDE{A}] 
  \item[\ABCDE{B}] 
  \item[\ABCDE{C}] 
  \item[\ABCDE{D}] 
  \item[\ABCDE{E}] 
  \item[\ABCDE{=}] 
\end{ABCDESlide}
}
{}
{}
{}

\Layout{0}
{{\TxSAMPLER}}
{
\begin{SAMPLERText}
  \item[\SAMPLER{S}] Es werden meist viele Medikamente genommen,
sowohl für die Herzinsuffizienz, als auch für die
Begleiterkrankungen. 
  \item[\SAMPLER{P}] Die Herzinsuffizienz ist als chronische
Erkrankung bekannt, entsprechende Befunde vorhanden. Oft
bestehen Erkrankungen oder gab es Ereignisse, welche 
den Herzmuskel geschädigt haben, wie
z.\,B. diverse Herzerkrankungen, frühere Herzinfarkte,
Bluthochdruck, u.\,a.
  \item[\SAMPLER{R}] Auslöser für akute Beschwerden können alle
Umstände sein, welche vom Herz eine ungewohnte Mehrarbeit
erfordern: Darunter fallen z.\,B. Infektionen und andere
Erkankungen, Stress und ein plötzlicher Blutdruckanstieg.
\end{SAMPLERText}
}
{
\begin{SAMPLERSlide}
  \item[\SAMPLER{S}] Viele
  \item[\SAMPLER{P}]  Herzinsuffizienz bereits bekannt,
  \\ Befunde, Bluthochdruck, Herzerkrnakungen, frühere
  Herzinfarkte
  \item[\SAMPLER{R}] Ungewohnte Belastung: Infekte,
  Erkrankungen, Blutdruckanstieg
\end{SAMPLERSlide}
}
{}
{}
{}


\MassnahmenMK{Spezielle
Maßnahmen}{Herzinsuffizienz,
symptomatisch}{m:herzinsuffizienz}{
\Cave{Patienten mit Herzinsuffizienz und ABCD-Problem oder 
Lungenödem mit brodelndem Atemgeräusch sind grundsätzlich notarztpflichtig. 

Auch bei kurzer Transportzeit hat die Stabilisierung des Patienten vor Ort Vorrang.

Bereits der Transport in das Fahrzeug kann gefährlich sein!}
}{}
% \MassnahmenNG{Spezielle
% Maßnahmen}{Herzinsuffizienz,
% symptomatisch}{\label{m:herzinsuffizienz}}{
% \begin{itemize}
%   \item Lagerung: Oberkörper hoch
%   \item Beengende Kleidung öffnen
%   \item Vitale Bedrohung beurteilen. \TxBeiVit \TxMassVit
%   \item \ce{O2}-Gabe gemäß \prettyref{m:sauerstoffberieselung}
%   \item Ursachenforschung
% \end{itemize}
% 
% \Cave{Kein Transport von Patienten mit hörbarem Lungenödem
% ohne Notarzt. Auch wenn das Spital \enquote{um{\textquotesingle}s Eck}
% ist.}
% }







\clearpage % TEMP temp clearpage
% \section{Koronare Herzkrankheit und Akutes Koronarsyndrom}
\section{Erkrankungen der Herzkrankzgefäße}
\index{Koronare Herzkrankheit}
\index{Herzkrankheit!Koronare}
\index{KHK}
\index{Akutes Koronarsyndrom}
\index{Koronarsyndrom!Akutes}
\label{topic:khk}
\label{topic:aks}
\label{topic:akutes-koronarsyndrom}

\Layout{0}{\TxBeschreibung}
{\DefinitionInPara{Die \T{Koronare Herzkrankheit} (\T{KHK}) ist eine
Erkrankung der \EE{Herzkranzgefäße} (Koronargefäße) bei der
es zu einer Verengung der Gefäße kommt.} Sie ist eine
\EE{chronische} Erkrankung, welche in akuten Notfallsituationen
resultieren kann. 

\DefinitionInPara{Das 
\T[Akutes Koronarsyndrom]{Akute Koronarsyndrom}
\index{Koronarsyndrom!akutes} (\I{ACS}\ZFootnote{\I{ACS}: \I{Acute Coonary Syndrome}})
ist ein  Symptomenkomplex, welcher
typischerweise bei Erkrankungen wie dem 
\E{Herzinfarkt} und der \E{stabilen} oder \E{instabilen Angina
pectoris} \EE{akut} auftritt.} 
Bei diesen Krankheitsbildern kommt es zu einer akuten \EE{Unterversorgung des
Herzmuskels} mit Blut und Sauerstoff aufgrund der
Durchblutungsstörung. Die Ursache sind in erster Linie
verengte oder verstopfte Herzkranzgefäße
(\I[Ischämie!Akutes Koronarsyndrom]{Ischämie},
\FullPrettyRef{topic:ischaemie}). Das akute Koronarsyndrom
ist häufig das Ergebnis einer vorbestehenden koronaren Herzkrankheit.
}
{%
\begin{itemize}
  \item Erkrankung der Herzkranzgefäße $\rightarrow$ Verengung
  \item Koronare Herzkrankheit (KHK, chronisch)
  \item Akutes Koronarsyndrom (ACS, akut)
  \begin{itemize}
    \item Symptomenkomplex, z.\,B. bei:
    \begin{itemize}
      \item Angina pectoris
      \item Herzinfarkt
    \end{itemize}
  \end{itemize}
\end{itemize}
}
{}
{}
{}

\Commentary[Akutes Koronarsyndrom]{
Das
\protect\T[Akutes Koronarsyndrom]{Akute Koronarsyndrom}
\protect\index{Koronarsyndrom!akutes}
ist ein Symptomenkomplex, welcher
typischerweise bei Erkrankungen wie dem
\E{Herzinfarkt} und der \E{stabilen} oder \E{instabilen Angina
pectoris} auftritt. Daneben gibt es noch andere
Erkrankungen, welche sich von den vorgenannten entweder
durch die Pathogenese (z.\,B. \protect\I{Prinzmetal-Angina}) oder
durch den Pathomechanismus (z.\,B. \protect\I{Tako-Tsubo-Kardiopathie})
wesentlich unterscheiden.
}
\Commentary[Koronare Herzkrankheit]{Eine koronare
Herzkrankheit ist oft anamnestisch bereits bekannt. Ein
Hinweis kann auch die Einnahme von ASS-Präparaten (z.\,B.
ThromboASS{\texttrademark}, HerzschutzASS{\texttrademark},
\ldots) oder ein Stent-Ausweis sein. 
Je nachdem welche Koronargefäße betroffen sind, spricht man
auch von einer 1-, 2- oder 3-Gefäßerkrankung (\Lang{engl.}
1-, 2-, 3-Vessel-Disease; bzw. \Lang{Abkz.} 1-, 2.-, 3.-VD).}

\begin{PictureSeriesWide}{}{Bilderserie:
\EEE{Koronargefäße}}
\PictureSeriesRowWide{3}
{./images/hirtler-aass/Herzvorderflaeche.pdf}
{Das Herz mit seinen Koronargefäßen \SourcePicture{Hirtler}}
{./images/wikipedia-cc-by/Hk_coronary_bionerd.jpg}
{Darstellung der
Herzkranzgefäße während einer Herzkatheteruntersuchung.
Dabei wird über die Leistenarterie ein Katheter
eingebracht und bis knapp vor das Herz zu den Abgängen der
Koronargefäße aus der Aorta vorgeschoben. Anschließend wird ein
Kontrastmittel appliziert um die Gefäße mittels Röntgen sichtbar zu machen.
\SourcePicture{WmCo
\enquote{Bionerd}, CC-BY-3.0}}
{./images/wikipedia-cc-by/Heart_coronary_artery_lesion-lq.jpg}
{Die Koronargefäße versorgen den
Herzmuskel von außen nach innen. \SourcePicture{Patrick J.
Lynch, CC-BY-2.5}} {./images/wikipedia-cc-by/Hk_coronary_bionerd.jpg}
{empty}
{}

\end{PictureSeriesWide}

\Layout{0}
{Akutes Koronarsyndrom}
{Da bei Behandlungsbeginn oft nicht klar ist, ob es sich um
einen Herzinfarkt oder einen Angina-pectoris-Anfall handelt,
wird bei Vorliegen der entsprechenden Symptome (Atemnot,
Thoraxschmerz, Todesangst, \ldots) einfach vom \T{Akuten
Koronarsyndrom}\ZFootnote{\Lang{Engl} Acute Coronary
Syndrome, \Abk{ACS}.} gesprochen. 

Es sind \E{zwei Auslösemechanismen} möglich:
Einerseits kann eine Minderdurchblutung des Herzmuskels
(Myokard)
  aufgrund einer oder mehrerer \E{Einengung(en) von
  Herzkranzgefäßen} (Koronararterien) die Ursache sein, oder
es kommt zu einem \E{erhöhtem Sauerstoffbedarf} des Herzens
bei
  Anstrengung.
In beiden Fällen kommt es zu einem \EE{Missverhältnis
zwischen Sauerstoffangebot und Sauerstoffbedarf} des
Herzens (\I[Ischämie!Angina pectoris]{Ischämie},
\FullPrettyRef{topic:ischaemie}), die Folge sind Schmerzen,
Atemnot und ein Leistungsverlust!


}
{}
{./images/hirtler-aass/Herzschmerz.pdf}
{Schmerzausdehnung beim ischämischen
Herzschmerz.}
{Hirtler.}

\Layout{27}
{Akutes Koronarsyndrom}
{}
{}
{./images/hirtler-aass/Herzschmerz.pdf}
{Typische Schmerzausdehnung bzw. -ausstrahlung beim ischämischen
Herzschmerz.}
{Hirtler.}



% \begin{PictureSeriesWide}{}{Bilderserie:
% \EEE{EKG-Veränderungen beim Herzinfarkt}}
% \PictureSeriesRowWide{2}
% {./images/wikipedia-pd/12_lead_generated_inferior_MI}
% {\Z{Typische Veränderungen im EKG bei einem
% Hinterwandinfarkt:
% Deutliche ST-Hebungen in den Ableitungen \EKG{II}, \EKG{III}
% und
% \EKG{aVF}} \SourcePicture{WMC/PD}}
% {./images/wikipedia-cc-by-sa/Heart_inferior_wall_infarct-00640}
% {Ischämisches Herzmuskelgewebe \ldots \SourcePicture{Patrick J. Lynch, medical illustrator; C. Carl Jaffe, MD, cardiologist; CC-BY 2.5}}
% {./images/wikipedia-cc-by-sa/Heart_inferior_wall_scar}
% {{\ldots} stirbt nach einger Zeit ab, es bildet sich eine (funktionslose) Narbe, die Herzleistung ist in Folge beeinträchtigt.\SourcePicture{Patrick J. Lynch, medical illustrator; C. Carl Jaffe, MD, cardiologist; CC-BY 2.5}}
% {empty}
% {}
% \end{PictureSeriesWide}


\Layout{0}
{\TxABCDE}
{~
\begin{ABCDEText}
  \item[\ABCDE{Szeneüberblick}]
  \item[\ABCDE{Eindruck}] 
  \item[\ABCDE{Bewusstsein}]
  \item[\ABCDE{Hauptbeschwerde}] 
  \item[\ABCDE{A}] 
  \item[\ABCDE{B}] 
  \item[\ABCDE{C}] 
  \item[\ABCDE{D}] 
  \item[\ABCDE{E}] 
  \item[\ABCDE{=}] 
\end{ABCDEText}
}
{
\begin{ABCDESlide}
  \item[\ABCDE{Sz}] 
  \item[\ABCDE{Ei}] 
  \item[\ABCDE{Be}] 
  \item[\ABCDE{Hb}] 
  \item[\ABCDE{A}] 
  \item[\ABCDE{B}] 
  \item[\ABCDE{C}] 
  \item[\ABCDE{D}] 
  \item[\ABCDE{E}] 
  \item[\ABCDE{=}] 
\end{ABCDESlide}
}
{}
{}
{}

 
\Layout{0}
{{\TxSAMPLER} beim akuten Koronarsyndrom}
{
\begin{description}
  \item[\SAMPLER{S}] Die Leitsymptome sind ein \EE{akuter,
  nicht atemabhängige, Brustschmerz} und \EE{Atemnot} (\E{Dyspnoe}).
  Der Schmerz ist meist hinter dem Brustbein lokalisiert und
  kann \EE{ausstrahlen}, häufig in den Kiefer, (linken)
  Arm und in den Rücken. Der Schmerz
  wird als brennend bzw. stechend beschrieben. Zusätzlich zu dem
  Schmerz verspürt der Patient auch oft ein \EE{Enge-}
  bzw. \EE{Druckgefühl} im Brustkorb, \Speech{\enquote{als
  würde jemand darauf stehen}}.
  Unter Umständen kann es auch zu einem \EE{Oberbauchschmerz} kommen,
  welcher als Verdauungsstörung oder
  Magen-Darm-Erkrankung fehlgedeutet werden kann.
  
  Bei Diabetikern kann es aufgrund der
  Nervenschädigungen auch zu einem schmerzlosen Verlauf
  kommen, man spricht dabei auch vom \enquote*{\T[stummer
  Infarkt]{stummen Infarkt}}\index{Infarkt, stummer}.
  Sonst steht normalerweise der \E{stechende/brennende}, nicht
  atemabhängige \II[Brustschmerz]{Brust-} oder
  \II{Oberbauchschmerz} zusammen mit \II{Atemnot} und
  \II{Todesangst} im Vordergrund. Der Brustschmerz kann in
  Richtung \E{Magen}, rechter oder linker \E{Hand}, \E{Kiefer}
  oder auch \E{Rücken} \EE{ausstrahlen}.
  
  \EE{Todesangst und Vernichtungsgefühl} sind
  häufig.
  Bei Auftreten eines \EE{kardiogenen Schocks} können die
  entsprechenden Schockzeichen beobachtet werden: Blässe,
  Kaltschweißigkeit, Hypotonie, etc.
  (\FullPrettyRef{topic:schock-kardiogener}).
  \item[\SAMPLER{M}] \II{Nitroglycerin} wird zur Erweiterung der
  Koronargefäßen von KHK-Patienten eingenommen. Dies kann
  entweder mittels eines Pflasters (Nitro-Pflaster), über
  Kapseln oder mittels eines Sprays (z.\,B.
  \I{Nitrolingual{\texttrademark}}), welcher bei Bedarf
  eingenommen wird, erfolgen. 
  Herzkranke Patienten nehmen daneben oft noch viele andere Medikemente ein.
  \item[\SAMPLER{P}]
  Bei den meisten Patienten ist die Angina
  Pectoris schon bekannt, bzw. es besteht eine \EE{koronare
  Herz-Krankheit} (\E{KHK}). Die KHK ist eine Erkrankung der
  Herzkranzgefäße, diese werden durch Ablagerung von Fetten
  und Kalk geschädigt und eingeengt, dadurch kann weniger Blut
  durchfließen. So kommt es zu einer Mangelversorgung des
  Herzmuskels und zu den Symptomen der Angina pectoris.
  \item[\SAMPLER{E}]
  Häufig wird über körperliche Anstrengung oder Aufregung
  vor Beginn des Anfalls berichtet.
\end{description}




}
{
\begin{itemize}
  \item[\SAMPLER{S}]
  \begin{itemize}
    \item Stechender/brennender \EE{Brustschmerz} mit
    Ausstrahlung oder \EE{Oberbauchschmerz}
    
    \enquote*{Stummer Verlauf} bei Diabetikern möglich
    \item \EE{Dyspnoe}
    \item Druckgefühl, \Speech{\enquote{als würde jemand darauf
    stehen}}
    \item Todesangst
    \item Ggfs. kardiogener Schock
  \end{itemize}
  \item[\SAMPLER{M}] Nitroglycerin: Pflaster, Kapseln, Spray
  \item[\SAMPLER{P}] Koronare Herzkrankheit (KHK)
  \item[\SAMPLER{E}] Körperliche Anstrengung
\end{itemize}
}
{./images/hirtler-aass/Herzschmerz.pdf}
{Schmerzausdehnung beim ischämischen
Herzschmerz.}
{Hirtler.}

\Layout{0}
{Angina pectoris}
{\index{Angina pectoris}
\label{topic:angina-pectoris}
\HeadingSub{\Lang{lat.} \enquote{Brustenge}}

\noindent\DefinitionInPara{Plötzliches Engegefühl in der Brust meist
mit Thoraxschmerz infolge einer \E{vorübergehenden}
Unter\-ver\-sorgung des Herzmuskels mit sauerstoffreichem Blut.}
Ein akuter Angina pectoris-Anfall ist primär \EE{nicht
sicher von einem Herzinfarkt unterscheidbar}. Die Symptome
treten eher bei Belastung auf und vergehen oft bei
Entspannung, das Vorübergehen der Beschwerden darf aber
natürlich nicht abgewartet werden!



Der akute Angina-Pectoris-Anfall gehört zu der Familie des
\EE{Akuten Koronarsyndroms}. Dazu gehört auch der
Herzinfarkt, dem der gleiche Mechanismus wie der Angina Pectoris zu
Grunde liegt, wobei es beim Infarkt im \E{Gegensatz} zur Angina
pectoris zu einem \E{totalen} Herzkranzgefäß\E{verschluss}
mit Absterben von Herzmuskelgewebe kommt.
\A{KOCH: nein; GABS: Doch. }}
{%
Missverhältnis zwischen Sauerstoffangebot und  
Sauerstoffbedarf
\begin{itemize}
  \item Einengung(en) von Herzkranzgefäßen
  \item Erhöhter Sauerstoffbedarf bei Anstrengung
  \item Symptome des akuten Koronarsyndrom,  eher bei
  Belastung, vergehen oft bei Entspannung
  \item Primär \EE{nicht sicher vom Herzinfarkt unterscheidbar}
\end{itemize}
}
{}
{}
{}




\Layout{0}
{Herzinfarkt}
{\HeadingSub{\Lang{Term.} \T{Myokardinfarkt}, \Lang{Abkz.} \T{MCI}.}
\index{Herzinfarkt}
\index{Infarkt!Herz-}
\index{Myokardinfarkt}
\index{MCI}
\index{Herzkranzgefäße!Verschluss}
\index{Koronargefäße!Verschluss}
\label{topic:herzinfarkt}
\label{topic:myocardinfarkt}
\label{topic:mci}

\noindent\DefinitionInPara{Verschluss eines oder mehrerer Herzkranzgefäße mit
  anschließendem Absterben des betroffenen Herzmuskelgewebes
  (\I[Ischämie!Herzinfarkt]{Ischämie},
\FullPrettyRef{topic:ischaemie}).}
Die Symptome entsprechen denen des akuten Koronarsyndroms, doch vergehen diese nicht mehr und es bleiben --
selbst bei optimaler Behandlung -- Schäden am Herzmuskel
bestehen. 


}
{%
\begin{itemize}
  \item \EE{Verschluss} von Herzkranzgefäßen
  \item \EE{Absterben} von Herzmuskelgewebe
  \item Symptome eines \EE{ACS}
\end{itemize}}
{}
{}
{}


\begin{PictureSeriesWide}{}{Bilderserie:
\EEE{Schädigung des Herzmuskel beim Herzinfarkt}}
\PictureSeriesRowWide{3}
{./images/wikipedia-cc-by-sa/Heart_normal_short_axis_section-00640}
{Das Herz im Querschnitt, gesehen von unten: Links der
rechte Ventrikel, rechts der muskelstarke linke Ventrikel.
 \SourcePicture{Patrick Lynch, CC-BY}}
{./images/wikipedia-cc-by-sa/Heart_inferior_wall_infarct-00640}
{Ischämisches Herzmuskelgewebe {\ldots}
\SourcePicture{Patrick J. Lynch, medical illustrator; C. Carl Jaffe, MD, cardiologist; CC-BY 2.5}} 
{./images/wikipedia-cc-by-sa/Heart_inferior_wall_scar}
{{\ldots} stirbt nach einger Zeit ab, es bildet sich eine
(funktionslose) Narbe, die Herzleistung ist in Folge beeinträchtigt. 
\SourcePicture{Patrick J. Lynch, medical illustrator; C. Carl Jaffe, MD, cardiologist; CC-BY 2.5}}
{empty}
{}
\end{PictureSeriesWide}

  
\Layout{0}{\TxKomplikationen}
{%
\begin{itemize}
  \item \EE{Herzrhythmusstörungen} (alle Arten), bis hin
  zum Kreislaufstillstand und Reanimation.
  Besonders häufig kommt es als Frühkomplikation zu plötzlichen, 
  lebensgefährlichen Rhythmusstörungen wie der ventrikulären 
  Tachykardie oder dem Kammerflimmern.
  \item Kardiogener Schock \index{Kardiogener
  Schock!beim Herzinfarkt}\index{Schock!Kardiogener!beim
  Herzinfarkt}
\end{itemize}
}{
\begin{itemize}
  \item \EE{Lebensgefahr!}
  \item Jede Herzrhythmusstörung als Komplikation möglich!
  \item Kammerflimmern als Frühkomplikation häufig!
\end{itemize}

}{}{}{}



\Layout{27}
{}
{}
{}
{./images/wikipedia-pd/12_lead_generated_ventricular_tachycardia}
{Lebensbedrohliche Rhythmusstörungen, wie hier die
ventrikuläre Tachykardie, sind häufige Komplikationen eines
Herzinfarktes.}
{WM/PD}




\MassnahmenMK{Spezielle Maßnahmen}{Thoraxschmerzen, jetzt beschwerdefrei}
{m:akutes-koronarsyndrom-beschwerdefrei}{}{}

\SlideCompensation{2}

\MassnahmenMK{Spezielle Maßnahmen}{Akutes Koronarsyndrom}
{m:akutes-koronarsyndrom}
{\Ca{Vitale Bedrohung!}

\Cave{Patienten mit einem akuten Koronarsyndrom sind grundsätzlich notarztpflichtig. 

Auch bei kurzer Transportzeit hat die Stabilisierung des Patienten vor Ort Vorrang. 

Der Notarzt ist auch für die Transportentscheidung wesentlich: 
Der Patient braucht u.\,U. \EE{ein entsprechend
ausgestattetes Spital} mit einem 
\EE{\E{dienst\-bereiten} Herzkatheterlabor}. Der Notarzt muss
entscheiden, ob das notwendig und sinnvoll ist.
}}{
}


\SlideCompensation{1}


\Layout{0}{Weitere Therapie}
{%
Herzinfarktpatienten werden \enquote*{besonders} behandelt
und normalerweise nicht auf eine beliebige Internistische
Station gebracht\citep{Kalla2006}. Der Notarzt muss
i.\,d.\,R. zwischen dem Einleiten einer Lyse-Therapie und
der Einweisung in ein Herzkatheterlabor entscheiden.

\begin{itemize}
  \item Bei der \II{Lyse-Therapie} versucht man mittels eines
  speziellen Medikaments eine Verstopfung aufzulösen.
  \item \index{PTCA}
  \index{Herzkatheterlabor}
  Mit einem \EE{Herzkatheter}
   versucht man die
  verengten oder verstopften Gefäße aufzudehnen. Dazu
  benötigt man ein Spital, welches über ein entsprechend
  ausgestattetes und
  \E{dienstbereites}
  \T{Herzkatheterlabor} verfügt.\opt{REGVIE}{\footnote{In
  Wien gibt es einen \E{Dienstplan} für die Katheterlabore,
  d.h. jedes Katheterlabor hat an bestimmten Tagen der Woche
  rund um die Uhr Bereitschaftsdienst. Eine Voranmeldung und
  vorangehende Absprache zwischen Notarzt und diensthabendem
  Oberarzt der Abteilung ist erforderlich. Die PTCA ist in
  Wien im AKH, Donauspital, KH Hietzing, Rudolfstiftung,
  Hanusch-Krankenhaus und Wilhelminenspital verfügbar.(Angaben
  ohne Gewähr und können sich
  jederzeit ändern.)}} \hfill\Commentary[Herzkatheter]{%
  Bei
  einer Herzkatheteruntersuchung (\I{Koronarangiographie})
  wird über die Leistenarterie ein Katheter eingebracht und 
  bis knapp vor das Herz zu den Abgängen der Koronargefäße 
  vorgeschoben und anschließend ein
Röntgen-Kontrastmittel appliziert. Dadurch kann mittels
Durchleuchtungsgeräten der Blutfluß in den  Koronargefäßen
beobachtet und beurteilt werden, sowie Verengungen und
Verschlüsse identifiziert werden (diagnostische
Herzkatheteruntersuchung).
Wird eine solche Stelle gefunden, kann diese oft auch gleich mittels des
Herzkatheters behandelt werden
(Intervention; \I{PCI} (Percutane Coronare
Intervention) oder \I{PTCA} (Perkutane Transluminale
Coronare Angioplastie)). Dabei wird ein Ballon
in die Engstelle eingebracht, welcher das Gefäß aufdehnt.
Anschließend wird ein zylindrisches Metallgeflecht, ein
sog.
\I{Stent}, eingebracht, welcher als Stütze für das Gefäß
dient und dieses längerfristig offen halten soll. Vgl.
\cite{Harrisons18DeKap230,Harrisons18DeKap246}
}
\end{itemize}

}
{
\begin{itemize}
  \item Lyse
  \item Herzkatheterlabor (PCI, PTCA)%
  \opt{REGVIE}{%
\begin{itemize}
  \item Dienstbereites Katheterlabor: eigener Dienstplan
\end{itemize}%
}  	
\end{itemize}

}{}{}{}


\begin{figure}[H]
\caption{EKG-Veränderungen beim Herzinfarkt:
Unterschiedliche Gefäße verursachen unterschiedliche
Veränderungen im EKG.} \subfloat[\Z{EKG-Veränderungen bei
einem Vorderwandinfarkt: ST-Hebungen in den
Brustwandableitungen.}
\SourcePicture{WM/PD}] {\includegraphics[width=\linewidth]{./images/wikipedia-pd/12_lead_generated_anterior_MI}}

\subfloat[\Z{EKG-Veränderungen bei einem Hinterwandinfarkt:
Deutliche ST-Hebungen (\enquote*{Katzenbuckel}) in den
Ableitungen \EKG{II}, \EKG{III} und
\EKG{aVF}}
\SourcePicture{WM/PD}] {\includegraphics[width=\linewidth]{./images/wikipedia-pd/12_lead_generated_inferior_MI}}
\end{figure}



\begin{Literary}
\fullcite{youtube:Herzinfarkt}

\fullcite{youtube:Herzinfarkt2}

\fullcite{Erc2010Sec5}

Sowie:
\cite{Harrisons18DeKap243,Harrisons18DeKap244,Harrisons18DeKap245}
\end{Literary}

\section{Herzrhythmusstörungen}
\label{topic:herzrhythmusstoerungen}
\index{Herzrhythmusstörungen}
\index{Rhythmusstörungen!Herz-}

Herzrhythmusstörungen sind häufig, allerdings auch oft sehr
unterschiedlich in ihrer Bedeutung: Manche sind chronisch
und eher wenig gefährlich, andere können plötzlich auftreten
und sogar akut lebensbedrohlich sein. Rhythmusstörungen
können sowohl selber das Problem darstellen, als auch \E{als
Komplikationen} im Rahmen anderer Krankheitsbilder (z.\,B.
Herzinfarkt, etc.) auftreten.

\Layout{0}{Arten von Herzrhythmusstörungen}
{%
Herzrhythmusstörungen können hervorgerufen werden durch eine
\EE{Störung der Frequenz} (zu langsam, zu schnell) oder
einer \EE{Störung der Regelmäßigkeit} (Rhythmus im engeren Sinn). 
Wenn man die
elektrischen Ströme mittels eines EKG darstellt, kann man
zusätzlich auch \EE{Störungen der Form der elektrischen
Herzströme} erkennen und beschreiben. Ein halbautomatischer
Defibrillator (AED) erledigt dies automatisch, er zeichnet
die elektrischen Herzströme auf und beurteilt diese.

Im Rahmen von Herzrhythmusstörungen kann es zu einer
\EE{Entkoppelung der elektrischen und der mechanischen
Herzaktion} kommen, z.\,B. wenn nicht jedem elektrischen
Impuls ein Herzschlag folgt.
}{
\begin{itemize}
  \item Störung der
  \begin{compactitem}
    \item Frequenz
    \item Regelmäßigkeit
    \item Form der elektrischen Herzströme
  \end{compactitem}
\end{itemize}

Entkoppelung von elektrischer und der mechanischer
Herzaktion möglich!
}{}{}{}

\SlideCompensationNoParskip{1}

\Layout{0}{Hier zählt die Geschwindigkeit: Tachy- und Bradykardie}
{%
%
\begin{itemize}
  \item \T{Bradykardie}:  Langsamer Herzschlag
  ({HF \textless} 60\,\nicefrac{}{\Min} {\tiny (beim Erwachsenen)})
  
  Leistungssportler können aufgrund ihres trainierten Körpers auch sehr
  niedrige Ruhepulsfrequenzen haben.
  \item \T{Tachykardie}:  Schneller
  \marginpar{{\footnotesize Tachometer im Auto:
  Geschwindigkeitsanzeiger, meistens zu schnell ;)}}
  Herzschlag
  (\Abk{HF} {\textgreater} 100\,\nicefrac{}{\Min} {\tiny
  (beim Erwachsenen)})
\end{itemize}

Bei der Beurteilung der Herzfrequenz muss auch immer an die Umstände 
(Ruhe, Belastung, \ldots) und an den
Erregungszustand des Patienten gedacht werden!
}{
\begin{itemize}
  \item Bradykardie \dotfill\ zu langsam
  \item Tachykardie \dotfill\ zu schnell
\end{itemize}
}{}{}{}


\begin{table}[H]
\FontTableBase
\begin{WidePar}
\caption{Übersicht: Wichtige und häufige
Herzrhythmusstörungen.\label{tab:herzrhythmusstoerungen}}

% Index-Einträge für Tabelle ---
\index{Bradykardie} 
\index{Tachykardie}
\index{Kammerflimmern}
\index{Flimmern!Kammer-}
\index{Vorhofflimmern}
\index{Flimmern!Vorhof-}
\index{Asystolie}
\index{Pulslose Ventrikuläre Tachykardie} 
\index{Tachylardie!Pulslose Ventrikuläre}
% --- Ende


\begin{minipage}{\linewidth}
\FontTableBase
\begin{tabulary}{\textwidth}{LLL}\addlinespace\toprule\addlinespace
\noindent\TabHeading{Rhythmus}
& 
\noindent\TabHeading{Bedeutung}
& 
\noindent\TabHeading{Bemerkung}
\\\addlinespace\midrule\addlinespace
\TabSub{Bradykardie}

{\scriptsize (HF\,{\textless}\,60\PerMin*)}
& 
Situations\-abhängig 
& 
\emph{\enquote{zu langsam}}
\\\addlinespace
\TabSub{Tachykardie} 

{\scriptsize (HF\,{\textgreater}\,100\PerMin*) }
& 
Situations\-abhängig 
& 
\emph{\enquote{zu schnell}}
\\\addlinespace\midrule\addlinespace
\TabSub{Kammerflimmern}
& Reanimation  & \emph{\enquote{schockbarer}} Rhythmus 
\\\addlinespace
\TabSub{Pulslose Ventrikuläre Tachy\-kardie} 
& Reanimation 
& \emph{\enquote{schockbarer}} Rhythmus 
\\\addlinespace
\TabSub{Asystolie} 

{\scriptsize (HF\,=\,0)} 
& 
Reanimation 
& 
Nulllinie, \emph{\enquote{nicht-schockbarer}} Rhythmus
\\\addlinespace
\TabSub{Vorhofflimmern}
& regellose Erregung im Vorhof

oft symptomlos 
& 
\E{Chronische Erkrankung.} Sehr viele Leute haben
Vorhofflimmern und leben ganz gut damit. 

Risikofaktor für:

\E{Tachykardie-Anfälle}

\E{Thrombenbildung}, dadurch erhöhtes Lungenembolie- und
Schlaganfallrisiko. Deshalb Vorbeugung mit
gerinnungs\-hemmenden (\enquote{blut\-ver\-dünnenden})
Medi\-ka\-menten:
z.\,B. \E{\enquote{Marcoumar}} ${\rightarrow}$ Patient ist
blutungsgefährdet
\\\addlinespace\midrule\addlinespace
\TabSub{Extrasystole}
&
Extraschläge
&
Vereinzelt bis \enquote{Salven} möglich

Von symptom-/bedeutungslos bis lebensbedrohlich
\\\addlinespace\bottomrule
\end{tabulary}
\end{minipage}
\end{WidePar}
\end{table}





\subsection{Besondere Rhythmen}





\Layout{0}{Reanimationspflichtige Rhythmusstörungen}
{~
\begin{itemize}
  \item \T{Pulslose Ventrikuläre Tachykardie} (\I{pVT}):
%   \index{Pulslose Ventrikuläre Tachykardie|TLike}
  \label{topic:pulslose-ventrikulaere-tachykardie}
  %
  Vorstufe zum Kammerflimmern. Das Herz kontrahiert
  so schnell, dass es sich nicht mehr füllen kann
  ${\rightarrow}$ \EE{keine Auswurfleistung}.
  \item \T{Kammerflimmern} (\I{AF}):
%   \index{Kammerflimmern|TLike}
  \label{topic:kammerflimmern}
  %
   Die elektrische
  Erregung im Herzen ist ungerichtet, als würden alle Autos
  eines Parkplatzes gleichzeitig in alle Richtungen losfahren
  ${\rightarrow}$ der Herzmuskel kontrahiert nicht geordnet,
  er \enquote{zittert} nur (wie ein Pudding), es wird \EE{keine
  Auswurfleistung} erreicht! Kammerflimmern darf nicht mit Vorhofflimmern verwechselt werden!
  \item \T{Asystolie}:
%   \index{Asystolie|TLike}
  \label{topic:asystolie}
  Im Herz entsteht keine elektrische
  Erregung, der Herzmuskel kontrahiert daher nicht und es erfolgt
  {keine Auswurfleistung}
  \item \T{Pulslose elektrische Aktivität} (\I{PEA}): Bei der
  pulslosen elektrischen Aktivität ist die Herzaktion von
  der ekeltrischen Aktivität \EE{entkoppelt}, d.\,h. das Herz
  reagiert nicht auf die Impulse des Reizleitungssystem. Der
  EKG-Befund kann zwar unauffällig sein, aufgrund der
  fehlenden Herzaktion entsteht trotzdem ein
  Kreislauffstillstand.
\end{itemize}


}
{
\begin{compactitem}
  \item Pulslose Ventrikuläre Tachykardie
  \begin{compactitem}
    \item Keine Zeit zum Füllen ${\rightarrow}$ kein Auswurf
  \end{compactitem}
  \item Kammerflimmern
  \begin{compactitem}  
    \item Wirre Erregung, keine sinnvolle Muskelarbeit ${\rightarrow}$ kein
    Auswurf.
  \end{compactitem}
  \item Asystolie
  \begin{compactitem}
    \item Keine elektrische Erregung, keine Muskelarbeit, kein
    Auswurf
  \end{compactitem}
  \item Pulslose elektrische Aktivität
  \begin{compactitem}
    \item Entkoppelung von Herzschlag und Impulsen des Reizleitungssystems
  \end{compactitem}
\end{compactitem}

}{}{}{}

\Layout{27}
{}
{}
{}
{./images/wikipedia-pd/12_lead_generated_sinus_rhythm}
{Unauffälliger EKG-Befund}
{WM/PD}

\begin{figure}[H]
\caption{Reanimationspflichtige Rhythmen}
\subfloat[\Z{Ventrikuläre Tachykardie} \SourcePicture{WM/PD}]
{\includegraphics[width=\linewidth]{./images/wikipedia-pd/Lead_II_rhythm_ventricular_tachycardia_Vtach_VT}}

\subfloat[\Z{Kammerflimmern}
\SourcePicture{WM/PD}]
{\includegraphics[width=\linewidth]{./images/wikipedia-pd/Lead_II_rhythm_generated_ventricular_fibrilation_VF}}

\subfloat[\Z{Asystolie}
\SourcePicture{WM/PD}]
{\includegraphics[width=\linewidth]{./images/wikipedia-pd/Lead_II_rhythm_generated_asystole}}
\end{figure}





\SlideCompensationNoParskip{1}
    
\Layout{0}{Vorhofflimmern}
{
\index{Vorhofflimmern|TLike}
\label{topic:vorhofflimmern}
Regellose Erregung im \E{Vorhof}, die oft symptomlos ist. 
(Die Herzkammern arbeiten dabei normal!) Vorhofflimmern kann
ständig bestehen oder auch nur manchmal (episodenhaft)
auftreten. Sehr viele Leute haben Vorhofflimmern und leben
ganz gut damit.
Häufige \EE{Komplikationen} sind jedoch:
\begin{compactitem}
  \item \EE{Tachykardie-Anfälle}
  \item \EE{Thrombenbildung} in den Vorhöfen: Dies stellt wiederum ein Risko dar für:
  \begin{compactitem}
    \item Arterielle Gefäßverschlüsse (\FullPrettyRef{topic:arterieller-gefaessverschluss})
    \item Schlaganfälle (\FullPrettyRef{topic:insult})
    \item Mesenterialinfarkte (\FullPrettyRef{topic:mesenterialinfarkt})
  \end{compactitem}
  Daher werden diesen Patienten meist \EE{gerinnungshemmende}
  (\enquote*{blutverdünnenden}) Medikamente zur Vorbeugung verschrieben (z.\,B. \I[Marcourmar!Vorhofflimmern]{Marcoumar}, der Patient ist blutungsgefährdet)
\end{compactitem}
}
{%
\begin{compactitem}
  \item Flimmern im \EE{Vorhof}
  \item Häufige und oft chronische Erkrankung
  \item Thrombenbildung, Tachykardieanfälle
  \item Oft \E{gerinnngshemmende Medikamente}
\end{compactitem}

}{}{}{}

\SlideCompensationNoParskip{1}


\Layout{0}{Extrasystolen}
{%
\index{Extrasystolen|TLike}
\label{topic:extrasystolen}
Extrasystolen sind \enquote{Extraschläge} des Herzens, sie können
vereinzelt oder gehäuft auftreten, eventuell auch
\enquote{salvenartig}. Je nach Häufigkeit haben Extrasystolen
unterschiedliche Auswirkungen auf den Kreislauf und können
symptom- bzw. bedeutungslos sein, oder auch Beschwerden
verursachen. Wenn die Extrasystolen die normale
Herzpumpfunktion massiv stören, können Sie auch
lebensbedrohlich sein.
}
{%
\begin{itemize}
  \item Extraschläge
  \item Vereinzelt od. gehäuft, evtl. \enquote{salvenartig}
  \item Unterschiedliche Auswirkungen auf Kreislauf: symptomlos bis
  lebensbedrohlich möglich
\end{itemize}

}{}{}{}





\subsection{Tachykarde Attacke}
\label{topic:tachykarde-attacke}


\Layout{0}{\TxBeschreibung}
{%
\DefinitionInPara{Plötzliche Tachykardie ohne erkennbare Ursache, wobei der Patient
(meistens) Symptome wahrnimmt und eine Herzfrequenz um bzw. über 160\PerMin*
vorliegt.} 
Tachykarde Attacken sind ein häufiges Notfallbild. Oft sind
chronische Erkrankungen die Ursache, diese
können jedoch sehr unterschiedlich sein und sind auch vom Fachmann bei weiterführenden Untersuchungen
oft nicht leicht zu ermitteln.
}{
\begin{itemize}
  \item HF\,$\geq$\,160\PerMin*
\end{itemize}
}{}{}{}

\Layout{0}
{\TxABCDE}
{
\begin{description}
%   \item[{\color{AltBlueLight}\usefont{T1}{PTSansNarrow-TLF}{m}{n}\selectfont A-Problem}] 
  \item[\ABCDE{B}] Dyspnoe kann in Folge der kardialen Mehrbelastung austreten
  \item[\ABCDE{C}] Es liegt klarerweise eine \EE{Tachykardie} vor. Der Puls
  ist i.\,d.\,R. hoch, meist über 160\PerMin* und evtl. nur schwach tastbar. Der Patient kann bei mangelnder Herzauswurfleistung kreislaufinsuffizient werden und einen kardiogenen Schock entwickeln. Im Extremfall ist auch ein Kreislaufstillstand möglich. 
%   \item[{\color{AltBlueLight}\usefont{T1}{PTSansNarrow-TLF}{m}{n}\selectfont D-Problem}]
%   \item[{\color{AltBlueLight}\usefont{T1}{PTSansNarrow-TLF}{m}{n}\selectfont E-Problem}]
\end{description}
}
{
\begin{itemize}
%   \item[{\color{AltBlueLight}\usefont{T1}{PTSansNarrow-TLF}{m}{n}\selectfont A-Problem}] 
  \item[\ABCDE{B}] Evtl. Dyspnoe
  \item[\ABCDE{C}] Tachykardie $\geq$\,160\PerMin*, Puls evtl. schwach tastbar; 
  
  Evtl. Kreislaufinsuffizienz, kardiogener Schock, Kreislaufstillstand
%   \item[{\color{AltBlueLight}\usefont{T1}{PTSansNarrow-TLF}{m}{n}\selectfont D-Problem}] 
%   \item[{\color{AltBlueLight}\usefont{T1}{PTSansNarrow-TLF}{m}{n}\selectfont E-Problem}]
\end{itemize}
}
{}
{}
{}

% \Layout{0}
% {\TxABCDE}
% {
% \begin{description}
%   \item[\ABCDE{A}] 
%   \item[\ABCDE{B}] 
%   \item[\ABCDE{C}] 
%   \item[\ABCDE{D}] 
%   \item[\ABCDE{E}] 
% \end{description}
% }
% {
% \begin{itemize}
%   \item[\ABCDE{A}] 
%   \item[\ABCDE{B}] 
%   \item[\ABCDE{C}] 
%   \item[\ABCDE{D}] 
%   \item[\ABCDE{E}] 
% \end{itemize}
% }
% {}
% {}
% {}
% 
% \Layout{0}
% {{\TxSAMPLE}}
% {
% \begin{description}
%   \item[\SAMPLER{S}] 
%   \item[\SAMPLER{M}] 
%   \item[\SAMPLER{P}]
%   \item[\SAMPLER{L}]
%   \item[\SAMPLER{E}]
% \end{description}
% }
% {
% \begin{itemize}
%   \item[\SAMPLER{S}]
%   \item[\SAMPLER{A}]
%   \item[\SAMPLER{M}] 
%   \item[\SAMPLER{P}] 
%   \item[\SAMPLER{L}]
%   \item[\SAMPLER{E}] 
% \end{itemize}
% }
% {}
% {}
% {}

\Layout{0}
{{\TxSAMPLE}}
{
\begin{description}
  \item[\SAMPLER{S}] Das Leitsymptom ist die \EE{Tachykardie} vor. 
  Der Patient
  klagt oft über ein \EE{\enquote{Herzklopfen}} und ist unruhig. Manchmal wird der Anfall von
  Angstgefühlen begleitet.
  
  Oft kommen \EE{Begleitsymtome} dazu, wie zum Beispiel Atemnot oder ein leichter
  Brustschmerz. Dies sind erste Zeichen, dass das Herz an
  die Grenze seiner Belastbarkeit stößt.
  \item[\SAMPLER{M}] Viele Patienten, bei denen
  Rhythmusstörungen bekannt sind nehmen entsprechende
  Medikamente \Z{(Antiarrhythmika)}. \Z{Bei manchen
  Rhythmusstörungen ist es üblich, dem Patienten bestimmte
  Medikamente zu verschreiben, welche er im Anfall nehmen
  soll um diesen zu beenden \Z{(\enquote*{pill in the
  pocket})}.}
  Das Erheben der tatsächlich eingenommenen sowie der verschriebenen Medikamente
  ist sehr wichtig, da diese Informationen des weiteren Behandlungsverlauf
  stark beeinflusst!
  \item[\SAMPLER{P}] Eine große Rolle spielt die \EE{Anamnese}: Da oft
  chronische Erkrankungen die Ursache sind, haben die
  Patienten diese Attacken immer wieder und kennen die
  Symptome bereits (und oft auch deren Behandlung). Oft sind die
  entsprechenden Grunderkrankungen schon seit langem bekannt und es 
  gibt bereits entsprechende \EE{Arztbriefe} oder \EE{Befunde}.
  \item[\SAMPLER{L}] Bei regelmäßig wiederkehrenden Beschwerden sind 
  Informationen über die vorherigen Episoden einzuholen.
  \item[\SAMPLER{E}] Die Ereignisse sind genau zu erfragen und zu
  dokumentieren, eventuell ergibt sich dadurch eine Spur zum auslösenden Faktor.
\end{description}
}
{
\begin{itemize}
  \item[\SAMPLER{S}]
  \begin{itemize}
    \item Tachykardie $\geq$\,160\PerMin*
    \item Spürbares Herzklopfen
    \item Evtl. Atemnot, leichter Brustschmerz, Angstgefühl
  \end{itemize}
  \item[\SAMPLER{M}] \Z{Antiarrhythmika.} Eingenommene und verschriebene Medikamente unbedingt sorgfältig erheben!
  \item[\SAMPLER{P}] Oft vorbekannte Erkrankung(en); Arztbriefe und Befunde!
  \item[\SAMPLER{L}] Frühere Episoden?
  \item[\SAMPLER{E}] Genau erfragen und dokumentieren!
\end{itemize}
}
{}
{}
{}



\Layout{0}{{\Lightning}Gefahr }
{%
Da das \EE{Herz plötzlich viel mehr Arbeit leisten muss},
kann es sein, dass es an seine \EE{Grenzen} stößt und
\EE{versagt} bzw. die Blutversorgung des Herzens für diese
Arbeit nicht mehr reicht und ein Angina-pectoris-Anfall oder
sogar ein Herzinfarkt\A{ KOCH: passt mit der Definition MCI
nicht überein, AP wäre besser oder? GABS: AP dazugenommen,
Infarkt aber gut möglich. } ausgelöst wird.
Weiters kann die Rhytmusstörung in eine
andere, akut lebensbedrohende Rhythmusstörungen umschlagen
und es dadurch zu einer akut lebensbedrohlichen Situation
kommen.
\Commentary[Tachykarde Attacke / Gefahren]
{Es gibt viele verschiedene Auslöser für tachykarde
Attacken. Manche Rhythmusstörungen sind tückisch, dass sie
von einem Moment umschlagen und einen Kreislaufstillstand
verursachen können. Zum Beispiel beim Kammerflattern besteht
oft eine \Conc{2}{1}-Überleitung vom Vorhof in die Kammern
mit einer resultierenden Frequenz um die 140\PerMin*. Kommt es
zu der gefürchteten \Conc{1}{1}-Überleitung wird die
Flatterfrequenz der Vorhöfe von \Range{250}{350}\PerMin* (!) direkt auf die
Kammern übergeleitet, dies kommt meist einem Herzstillstand
gleich. \cite{BLIM4}}

\bigskip
}{
\begin{itemize}
  \item Herz überfordert:
  \begin{compactitem}
    \item Angina pectoris
    \item Herzinfarkt
  \end{compactitem}
\end{itemize}
}{}{}{}





\MassnahmenMK{Spezielle Maßnahmen}{Tachykarde Attacke}
{m:tachykarde-attacke}
{}{}




\section{Störungen des Blutdrucks}
\subsection{Hypotonie}
\label{topic:hypotonie}

 
 
\Layout{0}{{\TxBeschreibung} und Richtwert}
{%
\DefinitionInPara{Als \T{Hypotonie} bezeichnet man einen 
unangemessen zu niedrigen Blutdruck.}
Beim Erwachsenen beträgt der \E{untere} Grenzwert des
systolischen Blutdrucks 100\MmHg*. 
Dieser Grenzwert ist jedoch sehr individuell. Viele
Menschen haben auch unter Normalbedingungen eine
systolischen Blutdruck von 90\MmHg* oder sogar noch
weniger, ohne Beschwerden zu haben. Dies muss bei der
Interpretation des Wertes berücksichtigt werden. Bei Kindern
gelten grundsätzlich andere Nomalwerte.
\citep{PraxisleitAllgemeinmed6} }{%
\begin{itemize}
  \item Erwachsener: RR\textsubscript{sys} {\textless}
  ca. 100\MmHg*
  \item Individuell unterschiedlich
  \item Kinder haben andere Normalwerte
\end{itemize}
}{}{}{}


\subsubsection{Kollaps, Synkope}
\label{topic:kollaps}
\label{topic:synkope}
\index{Kollaps|TLike}
\index{Synkope|TLike}


\Layout{2}{{\TxBeschreibung} und Ursachen}
{%
\DefinitionInPara{Als Kollaps bzw. Synkope bezeichnet man 
eine kurzfristige Kreislaufregulationsstörung.} 
Manche Autoren unterscheiden zwischen beiden Begriffen und 
ordnen der Synkope eher kardiale Ursachen zu. 

Ein Kollaps kann seine Ursachen in verschiedenen
Organssystemen haben:
\begin{itemize}
  \item Herz (z.\,B. Herzrhythmusstörungen \Z{(\I{Adams-Stokes-Anfall})})
  \item Gefäßen (z.\,B. Hitzebedingte Erweiterung
  der Gefäße mit \enquote*{Versacken} des Blutes in den
  Beinen)
  \item Blut (z.\,B. Blutarmut (\I{Anämie}))
  \item Hirn
  \item Hirngefäßen (z.\,B. TIA,
  \prettyref{topic:tia} ) \ldots
  \item Sonstige Ursachen und Faktoren: Hypoglykämie, 
  Menstruation
\end{itemize}

Eine Abklärung ist präklinisch oft nicht möglich, und auch
unter stationären Bedingungen oft langwierig und nicht
selten ergebnislos. Da jedoch schwerwiegende und
lebensbedrohliche Diagnosen möglich sind, muss bei unklaren
Kollapszuständen eine gründliche Abklärung erfolgen.

Es gibt jedoch auch typische \enquote*{Kollaps-Situationen}:
Ein schwüler Tag in einer überfüllten U-Bahn-Garnitur, nicht
gefrühstückt, auf dem Weg zu einem Vorstellungstermin und
seit gestern in der Menstruation. Oder im stickigen
Veranstaltungszelt bei 40\Deg , wobei der der ganze
Flüssigkeitskonsum aus 5 (entwässernden) Dosen
RedBull{\texttrademark}
bestand. So mancher war auch ob eines Begräbnisses so
ergriffen, dass er sich fast in die Grube neben den Sarg
dazugelegt hätte \ldots 
}{%

}{}{}{}


\Layout{0}{\TxSymptome}
{%
Der Patient berichtet über Schwindel, evtl. ist er
zusammengesackt, war kurz bewusstlos. Der Patient wacht oft
auf dem Boden liegend auf und weiß nicht, was geschehen ist.
Er erzählt  es wäre ihm schwarz vor den Augen geworden. Der
Blutdruck ist meist niedrig, der Patient ist eventuell
bradykard oder tachykard.\bigskip
}{
\begin{compactitem}
  \item Schwindel
  \item Evtl. kurze Ohnmacht, Schwarzwerden vor Augen
  \item RR ${\downarrow}$, evtl. Bradykardie oder Tachykardie
\end{compactitem}
}{}{}{}





\MassnahmenMK{Spezielle
Maßnahmen}{Kollaps/Synkope}{m:kollaps}{\EE{Taktik}:
\Speech{Ausführliche Untersuchung und symptomatische Behandlung 
\Ca{Verwechslung mit lebensbedrohlichen
Differentialdiagnosen möglich!}}}{}

% \MassnahmenNG{Spezielle
% Maßnahmen}{Kollaps/Synkope}{\label{m:kollaps}}{ \EE{Taktik}:
% \Speech{Ausführliche Untersuchung und symptomatische Behandlung 
% \Ca{Verwechslung mit lebensbedrohlichen Differentialdiagnosen möglich!}}
% 
% \begin{itemize}
%   \item Vitale Bedrohung einschätzen, auf \EE{Differentialdiagnosen} untersuchen, z.\,B.:
%   \begin{compactitem}
%     \item Insult
%     \item Herzrhytmusstörungen
%     \item Akutes Koronarsyndrom (evtl. schmerzlos!)
%     \item Krampfanfall
%     \item Schädel-Hirn-Trauma (SHT)
%     \item Zuckerstoffwechselstörung (Hyper-, Hypoglykämie)
%     \item Exsikkose
%     \item \ldots
%   \end{compactitem}
%   Es sind dabei alle zur Verfügung stehenden
%   Möglichkeiten auszuschöpfen, z.\,B.:
%   \begin{compactitem}
%     \item Traumacheck (Sturz?),
%     \item Neurocheck inkl. Blutzuckermessung,
%     \item Temperatur,
%     \item Ausführliche Anamnese bzw. Fremdanamnese 
%     \item \ldots
%   \end{compactitem}
%   \item \EE{Ursachenforschung} (Begleiterkrankungen (Infekt, \ldots), letzte Mahlzeit, Hitze, \ldots)
%   \item Lagerung: Nach Ausschluss von Herzbeschwerden,
%   Atemnot und relevanten Verletzungen: Beine hoch
%   \item Transportentscheidung: Hospitalisierung zur Abklärung
%   anstreben\newline
%   Abteilung: Innere Medizin
% \end{itemize}
% 
% }

\subsection{Hypertonie}
\label{topic:hypertonie}
Die \T{Hypertonie}
(\T{Bluthochdruck})ist eine meist
symptomlose \textit{chronische} Erkrankung welche auf Dauer
schwere Schäden am Körper anrichten kann und im Zuge von
plötzlichen Blutdruckkrisen auch akut Probleme bereiten
kann.

\begin{itemize}
  \item Chronische Hypertonie: Obergrenze derzeit bei \RRmmHg{140}{90} (beim
  Erwachsenen).
\end{itemize}

\Cave{Bei \EE{Schwangeren} sind die Grenzwerte genau zu beachten!
(\EE{EPH-Gestose}, \prettyref{topic:eph-gestose})!}

\Layout{0}{Ursachen}
{
\begin{itemize}
  \item Widerstandserhöhung (durch Engstellung der Gefäße)
  \item Vermehrte Herzarbeit
  \item Sympathikusaktivierung, Adrenalin und Cortisol
  (Stresshormone) durch Dauerstress,  Bewegungsmangel, genetische Veranlagung
  \item u.\,a. \ldots
\end{itemize}

}{

}{}{}{}


\Layout{60}{Langzeitfolgen}{%
Die Hypertonie hat in der Notfallmedizin auch eine wichtige
indirekte Bedeutung: Sie ist Ursache für viele andere
Erkrankungen, die ihrerseits akut lebensbedrohlich werden
können.
Zwei Wirkungen der Hypertonie sind dabei besonders wichtig:
Die \EE{Gefäßschädigung} und die \EE{Belastung des Herzens
durch den hohen Gegendruck}.

Bei der Gefäßschädigung kommt es zur Ablagerung durch Kalk
und sonstige Plaques und über längere Zeit zur \E{Verengung
oder Verstopfung der Gefäße}. Je nach Organ kann das sehr
massive Auswirkungen haben. Im Herz führt es zur \E{Angina
Pectoris} oder zum \E{Herzinfarkt}, im Hirn zu
\E{Schlaganfällen}, in den Nieren zu einer
\E{Niereninsuffizienz}, usw. Selbst große Gefäße wie die
Hauptschlagader (Aorta) können derartig geschädigt werden,
dass sie reißen können (Aortenaneurysma) -- eine für den
Patienten sehr lebensgefährliche Komplikation.

Die Schädigung des Herzmuskels ist die Folge von dem
erhöhten Druck gegen den das Herz pumpen muss. Der Muskel
vergrößert sich um mit der Belastung fertig zu werden, kann
aber dann nicht mehr ausreichend mit sauerstoffreichen Blut
versorgt werden.

Die Hypertonie ist somit ein wichtiger
\E{Risikofaktor} für viele schwere Notfälle. Daher ist
es auch verständlich, dass die Behandlung der Hypertonie
einen großen Stellenwert in der medizinischen Versorgung der
Bevölkerung hat.
}{
\begin{itemize}
  \item \EE{Gefäßschädigung}
  \begin{itemize}
    \item Herzkranzgefäße ${\rightarrow}$ Infarkt
    \item Hirngefäße ${\rightarrow}$ Schlaganfall
    \item Nierengefäße ${\rightarrow}$ Niereninsuffizienz ${\rightarrow}$
    Dialyse ${\rightarrow}$ typischer KTW-Patient
    \item Hauptschlagader ${\rightarrow}$ Aneurysma
  \end{itemize}
  \item \EE{Herzmuskelschädigung} \par  (Herz muss gegen
  größeren Widerstand pumpen)
  \begin{itemize}
    \item Krankhafte Vergrößerung des Herzmuskels
  \end{itemize}
\end{itemize}
}
{./images/teaser/imgp0481.jpg}
{}
{GABS}




\subsubsection{Hypertensive Krise und Hypertensiver Notfall}
\label{topic:hypertensive-krise}
\label{topic:hypertensiver-notfall}

\Layout{2}{\TxBeschreibung}
{%
Eine plötzliche, sehr starke Erhöhung des Blutdrucks wird
als hypertensive Krise oder hypertensiver Notfall
bezeichnet. Der Unterschied zwischen den beiden Diagnosen
ist das Fehlen bzw. Vorhandensein von typischen Symptomen,
s.
\prettyref{tab:hxpertension-krise-notfall}.

}{}{}{}{}

\Layout{0}{\TxSymptome}
{%
Vgl. \FullPrettyRef{tab:hxpertension-krise-notfall}.
Ein erhöhter Blutdruck alleine ist oft symptomlos und wird
vom betreffenden Patienten oft gar nicht bemerkt.
Klassisch kann er sich durch \I{Nasenbluten}
(\T{Epistaxis}\index{Epistaxis}) bemerkbar machen, welches
ein typisches Symptom für einen plötzlich erhöhten Blutdruck
ist.

Wenn Symptome wie Kopfschmerzen, Schwindel und Übelkeit,
Brustschmerzen, Sehstörungen, etc. auftreten, ist dies ein
Zeichen für die blutdruckbedingte Störung von Organen
(\E{Organstörungen}). Diese Symptome einer Organstörung sind
typische Leitsymptome für einen \E{hypertensiven
\EE{Notfall}}.

}{
\FullPrettyRef{tab:hxpertension-krise-notfall}.
}{}{}{}

\Layout{57}{}{}{}{
\begin{tabularx}{\textwidth}{XX}
\toprule\addlinespace
\noindent\TabHeading{Hypertensive Krise}
&
\noindent\TabHeading{Hypertensiver Notfall}
\\\addlinespace\midrule\addlinespace
\Abk{RR}~{\textgreater}~\RRmmHg{230}{130} (Richtwert) 
bzw. deutliche individuelle Erhöhung

\EE{Patient fühlt sich gut}, evtl. Nasenbluten
 &
\Abk{RR}~{\textgreater}~\RRmmHg{230}{130} (Richtwert)
bzw. deutliche individuelle Erhöhung

\EE{+ Symptome einer Organstörung}
\\\addlinespace
Zufallsbefund?

Oder steht der Blutdruck doch in Zusammenhang mit der
Berufung?

\textit{Patient hat überhaupt keine Beschwerden?} Warum hat
er Sie dann überhaupt gerufen?

Wenn es wirklich nur ein Zufallsbefund ist
(\enquote*{Gesundheitstage},
\enquote*{Vorsorgeuntersuchung}):
&
\begin{itemize}
    \item Kopfschmerz
    \item Sehstörungen
    \item Brustschmerz
    \item Übelkeit, Erbrechen, Schwindel
\end{itemize}
\\\addlinespace
Maßnahmen: Eine notfallmäßige Blutdrucksenkung ist (eher) nicht erforderlich.
&
Vitale Bedrohung: \Ca{Gefahr des Herzversagens: Das Herz muss
plötzlich gegen großen Druck pumpen}

Eine frühzeitige Blutdrucksenkung 
ist notwendig. 

Evtl. Anwendung der Notkompetenz gem. SanG. 
\\\addlinespace\bottomrule
\end{tabularx} 
}
{Unterscheidung Hypertensive Krise und Hypertensiver Notfall
\protect\citep{WHOISH2003,JNC7Report,BLIM3,IntroClinEmergMed4,Vlcek2008}.\label{tab:hxpertension-krise-notfall}}
{}

\Layout{0}
{\TxABCDE} {
\begin{description}
  \item[\ABCDE{Ersteindruck}] Oft hat der Patient einen hochroten Kopf. 
  \item[\ABCDE{Hauptbeschwerde}] Häufig ist Nasenbluten das
  erste Zeichen einer akuten Hypertonie, daneben kommt es
  oft zu neurologischen Symptomen wie Kopfschmerzen und
  Sehstörungen. Evtl. kommt es zu Brustschmerzen, Übelkeit,
  Erbrechen und Schwindel.
  \item[\ABCDE{B}] Evtl. kommt es aufgrund der kardialen Belastung zur Atemnot.
  \item[\ABCDE{C}] In sehr schweren Fällen, v.\,a. wenn eine
  kardiale Grunderkrankung besteht, können sich Symptome
  eines kardiogenen Schocks zeigen.
  \item[\ABCDE{D}] Es ist eine deutliche Blutdruckerhöhung
  (Richtwert \Abk{RR}~{\textgreater}~\RRmmHg{230}{130}
  bzw. deutliche individuelle Erhöhung) festzustellen
\end{description}
}
{
\begin{itemize}
  \item[\ABCDE{EE}] Hochroter Kopf 
  \item[\ABCDE{HB}] Nasenbluten, Kopfschmerzen, Sehstörungen, Brustschmerzen, Übelkeit, Erbrechen, Schwindel
  \item[\ABCDE{B}] Evtl. Dyspnoe
  \item[\ABCDE{C}] Evtl. kardiogener Shock
  \item[\ABCDE{D}] RR $\uparrow$ (Richtwert\Abk{RR}~{\textgreater}~\RRmmHg{230}{130} 
  bzw. deutliche individuelle Erhöhung)
\end{itemize}
}
{}
{}
{}

\Layout{0}
{{\TxSAMPLE}}
{
\begin{description}
  \item[\SAMPLE{S}] 
  \item[\SAMPLE{M}] 
  \item[\SAMPLE{P}]
  \item[\SAMPLE{L}]
  \item[\SAMPLE{E}]
\end{description}
}
{
\begin{itemize}
  \item[\SAMPLE{S}]
  \item[\SAMPLE{A}]
  \item[\SAMPLE{M}] 
  \item[\SAMPLE{P}] 
  \item[\SAMPLE{L}]
  \item[\SAMPLE{E}] 
\end{itemize}
}
{}
{}
{}

\Commentary[Hypertensive Krise/Notfall]{(Die Unterscheidung
zwischen \enquote{Hypertensiver Krise} und \enquote{Hypertensiven Notfall} ist in der
Literatur nicht einheitlich. Oft wird auch nur
zwischen \enquote{Hypertensiver Krise ohne Symptome} bzw. \enquote{mit
Symptomen} unterschieden oder generell eine andere
Einteilung gewählt. Von diesen Spitzfindigkeiten abgesehen
gilt das Vorhandensein von Symptomen durchgehend als
entscheidendes Kriterium.
\citep{ClassenInnere5,KlinLeitAlle3,Fauci2008})}
% Renz-Polster2006 p442
% IntroClinEmergMed4 Kap. 28, p 393

\MassnahmenMK{Spezielle Maßnahmen}{Hypertensive
Krise}{m:hypertensive-krise} {%
\EE{Taktik}: \Speech{Symptomatische Therapie}}{}

% \MassnahmenNG{Spezielle Maßnahmen}{Hypertensive Krise}{\label{m:hypertensive-krise}}
% {%
% \EE{Taktik}: \Speech{Symptomatische Therapie}
% 
% \begin{itemize}
%   \item Lagerung: Oberkörper hoch
%   \\
%   Bei Nasenbluten: Kopf nach vorne halten.
%   \item Beengende Kleidung öffnen
%   \item \ce{O2}-Gabe gemäß \prettyref{m:sauerstoffberieselung}
%   \item Transportentscheidung: Abt. f. Innere Medizin; 
%   
%   bei Nasenbluten: HNO im Hintergrund
% \end{itemize}
% Eine notfallmäßige Blutdrucksenkung ist \E{nicht}
% erforderlich. }

\MassnahmenMK{Spezielle Maßnahmen}{Hypertensiver
Notfall}{m:hypertensiver-notfall} {
\EE{Taktik}: \Speech{\Ca{Vitale Bedrohung!}
Rasche Blutdrucksenkung (veranlassen)}}{}

% \MassnahmenNG{Spezielle Maßnahmen}{Hypertensiver Notfall}{\label{m:hypertensiver-notfall}}
% {
% \EE{Taktik}: \Speech{\Ca{Vitale Bedrohung!}
% Rasche Blutdrucksenkung (veranlassen)}
% 
% \begin{itemize}
%     \item Lagerung: Oberkörper hoch
%     \\
%     Bei Nasenbluten: Kopf nach vorne halten.
%   \item Beengende Kleidung öffnen
%   \item \TxMassVit
% \end{itemize}
% 
% \Customized{
% \begin{description}
%   \item[\CustAsbLvWien, \CustAsbFlo] \Abk{NFS} dürfen lt.
% Algorithmus Nitro-Spray verabreichen. Ist die Gabe
% erfolgreich, kann u.\,U. von der Anforderung eines Notarztes
% abgesehen werden.
% \end{description}
% }
% }


\bigskip

\section{Gefäßaussackungen: Aneurysmen}
\label{topic:aneurysma}

\subsection{Bauchaortenaneurysma}
\label{topic:bauchaortenaneurysma}
\nocite{Meron2004}

%%%%%%%%%%%%%%%%%%%%%%%%%%%%%%%%%%%%%%%%%%%%%%%%%%%%%%%%%%%
% TODO : Aneurysma Überarbeitung
\TODO{GABS}{????-??-??}{Ticket 34: Erklärung Aneurysma.}{Dieser Teil ist nicht gut formuliert und/oder unvollständig, 
   er benötigt eine Überarbeitung!}
%%%%%%%%%%%%%%%%%%%%%%%%%%%%%%%%%%%%%%%%%%%%%%%%%%%%%%%%%%%
\Layout{0}{\TxBeschreibung{} und Ursachen}
{%%% Gerhard Gabriel <gerhard.gabriel@grissu.org>
Bei der Entwicklung eines Bauchaortenaneurysmas ist zumeist
eine krankhafte Veränderung der Aortenwand
verantwortlich. Als Folge eines Unfallgeschehens,
oder auch nach Hebeversuchen schwerer Lasten, kann die
Gefäßwand im Bereich eines Aneurysmas einreißen. 
}
{\SymbolText}{}{}{}

\Layout{0}{\TxSymptome}{%
Je nach
Größe der Schädigung und entsprechendem Schweregrad der Blutung
im Läsionsbereich, kommt es binnen kürzester Zeit unter
heftigem, akutem Bauchschmerz zum Tod. Kleinere Läsionen
sind ebenfalls von heftigem, akut einsetzendem Bauchschmerz
begleitet. Es bildet sich ein Blutungsschock (hypovolämischer
Schock) aus, der unbehandelt zum Tode führt.
}{
\begin{itemize}
  \item Heftigste Schmerzen
  (\enquote{Vernichtungsschmerz})
  \item Hoher Blutverlust ${\rightarrow}$ schnell im
  Schock!
  \item Evtl. bereits bekanntes Aneurysma
  ${\rightarrow}$\,Anamnese
  \item Evtl. tastbare pulsierende Schwellung
\end{itemize}
}{}{}{}



Oft ist bei den Patienten das Aneurysma bereits bekannt und
war noch nicht groß genug um behandelt zu werden
(${\rightarrow}$\,Anamnese!). Ansonsten ist es schwer das
Aneurysma präklinisch zu diagnostizieren. Im Vordergrund
steht der Volumenverlust und der entstehende hypovolämische
Schock. 

\bigskip

% \clearpage
\section{Gefäßverschlüsse in den Extremitäten}
\label{topic:gefaessverschluesse-extremitaeten}

\Layout{0}{Unterscheidung}
{%
Grundsätzlich unterscheidet man je nach betroffenem Gefäß
zwischen \EE{arteriellen} und \EE{venösen}
Gefäßverschlüssen. 
Beim \E{arteriellen Verschluss} besteht
das Hauptproblem in der Unterversorgung der Extremität mit
Blut. Beim \E{venösen Verschluss} kommt es zu einer Ab- und
Rückflussbehinderung. Zusätzlich besteht die Gefahr, dass sich
ein Gerinnsel losreißt und in einem anderen Teil des
Körpers Schaden anrichtet\footnote{Zum Beispiel in der
Lunge im Rahmen einer Lungenembolie}.
}{
\begin{itemize}
  \item \textbf{\textcolor[rgb]{0.7882353,0.0,0.08627451}{Arteriell}}
  \item \textbf{\textcolor{blue}{Venös}}:
  \index{tVT}\T{tVT} (tiefe Venen Thrombose)
\end{itemize}
}{}{}{}



\subsection{Arterieller peripherer Gefäßverschluß}
\label{topic:arterieller-gefaessverschluss}

\Layout{0}{Beschreibung}
{%
Beim \E{arteriellen Gefäßverschluss} wird das
zuführende Gefäß verschlossen. Es kommt zu einer schmerzhaften 
Unterversorgung der Extremität mit Blut
(\I[Ischämie!Arterieller peripherer
Gefäßverschluß]{Ischämie},
\FullPrettyRef{topic:ischaemie}). Der komplette Verschluss passiert plötzlich. }
{
\begin{itemize}
    \item Verschluß eines blut\E{zuführenden} Gefäßes
\end{itemize}
}{}{}{}

\Layout{2}{Ursachen}
{%
Meistens liegt eine 
chronische Schädigung der
Gefäße infolge anderer Erkankungen oder aufgrund eines
ungesunden Lebensstils vor. Ursachen für solch eine
Gefäßschädigung sind z.\,B. die Zuckerkrankheit (\T{Diabetes
mellitus}, s.\ref{topic:diabetes}), Hypertonie, erhöhte
Blutfette oder das Rauchen. Die Gefäße verkalken und
verändern sich auch sonst nicht zu ihrem Vorteil.
}
{
\begin{itemize}
  \item Chronische Gefäßschädigung, z.\,B. Bei Diabetes mellitus, Hypertonie \ldots
\end{itemize}
}{}{}{}



\Layout{0}{Chronische Grunderkrankung: pAVK}{
Wenn die Gefäße nur verengt sind, spricht man von
der \E{peripheren
arteriellen Verschluss\-krankheit}, abgekürzt \T{paVK}. Sie ist eine
häufige chronische Erkrankung. (\E{Achtung}: Obwohl es
\enquote*{\emph{Verschluss}krankheit} heißt, sind die Gefäße
nicht komplett verschlossen, sondern nur verengt!)
}{
\begin{itemize}
    \item Periphere arterielle Ver\-schluss\-krank\-heit
\end{itemize}

}{}{}{}

 

\Layout{}{}{}{}{}{}{}

\Layout{0}{\TxSymptome}
{%
Die Symptome ergeben sich aus der Unterversorgung der
Extremität mit sauerstoffreichem Blut: sie ist blass,
kalt, tut höllisch weh und die Rekap-Zeit ist nicht
mehr sinnvoll messbar. Zusätzlich wird die Extremität
gefühllos und gelähmt. Der Patient kann ein Kribbeln
empfinden (\E{\enquote{Ameisenlaufen}}).\ZFootnote{Im
englischen werden die Symptome als die
\Speech{\enquote{five P's of limb ischemia}}
zusammengefasst: {Pain} (schmerzend)
{Pale} (blaß),
{Pulseless}  (pulslos),
{Paresthesia} (gefühllos),
{Paralysis}} (Lähmung).

% \Captionof{table}{Die 5 Zeichen des arteriellen Gefäßverschluss.\label{tab:5p}}
% 
% \begin{tabulary}{\linewidth}{CLL}
% \midrule\addlinespace
% \centering 1. & \noindent\Z{Pain}        & \noindent\EE{schmerzend}
% \\
% \centering 2. & \noindent\Z{Pale}        & \noindent\EE{blaß}
% \\
% \centering 3. & \noindent\Z{Pulseless}    & \noindent\EE{pulslos}
% \\
% \centering 4. & \noindent\Z{Paresthesia} & \noindent\EE{gefühllos}
% \\
% \centering 5. & \noindent\Z{Paralysis}   & \noindent\EE{gelähmt}
% \\\addlinespace\bottomrule
% \end{tabulary}
}{%
\begin{compactitem}
  \item Schmerz
  \item Blässe
  \item Kein Puls
  \item Gefühllosigkeit
  \item Lähmung
\end{compactitem}
}{}{}{}


\MassnahmenMK{Spezielle Maßnahmen}{Arterieller
Gefäßverschluss}{m:arterieller-gefaessverschluss}{
\EE{Taktik}:
\Speech{Symptomatische Therapie und zügiger Transport an
geeignete Einrichtung}}{}

% \MassnahmenNG{Spezielle Maßnahmen}{Arterieller
% Gefäßverschluss}{\label{m:arterieller-gefaessverschluss}}{
% \EE{Taktik}:
% \Speech{Symptomatische Therapie und zügiger Transport an geeignete Einrichtung}
% 
% \begin{itemize}
%   \item Lagerung: Extremität hängen lassen, weich und warm lagern
%   \item Schonend transportieren
%   \item Nüchtern
%   \item Transportentscheidung: Abteilung für Chirurgie
%   (Gefäßchirurgie)
% \end{itemize}
% }

\subsection{Tiefe Venenthrombose}
\label{topic:venenthrombose}
\label{topic:thrombose}
\index{Thrombose}
\index{Venenthrombose!tiefe}
\index{Beinvenenthombose}
\index{tVT}

\Definition{Verstopfung einer Vene durch ein Blutgerinnsel
(Thrombus\footnote{Thrombus: Blutgerinnsel. Im Gegensatz
dazu \T{Embolus}: Losgelöster Thrombus, welcher woanders ein
Gefäß verstopft und eine \T{Embolie} verursacht.}).}


\Layout{0}{Ursachen und Risikofaktoren}
{%
Ursächlich für die Bildung eines Blutgerinnsels ist meistens
eine verminderte Blutflussgeschwindigkeit, welche wiederum
verschiedene Ursachen haben kann. 
Die \EE{venöse Insuffizienz} ist ein weit verbreitetes
chronisches Leiden.
 Es kommt  zu einem Rückfluss und Aussackungen von
Venen, welche im fortgeschrittenen Stadium als hässliche
\E{Krampfadern} deutlich sichtbar sind. Schuld sind oft
insuffiziente Venenklappen, stehende bzw. sitzende
Tätigkeiten und fehlende Bewegung.

\index{Economy class Syndrom}
Eine länger dauernde \EE{Immobilisation} oder
\EE{Bettlägrigkeit} sind auch typische Ursachen für einen
verminderten venösen Blutfluss. Darunter fallen auch lange
Reisen ohne Bewegung: So werden Thrombosen aufgrund von
langen Flugreisen auch scherzhaft als \T{Economy class
Syndrom} bezeichnet.
Erschwerend wirkt sich auch eine \EE{unzureichende
Flüssigkeitsaufnahme} aus, da dann der flüssige Anteil des
Blutes vermindert sein kann und das Blut folglich
dickflüssiger\ZFootnote{visköser} wird.

Ein besonderer Risikofaktor ist der Zigarettenkonsum. Die
Einnahme von hormonellen Verhütungsmitteln (\EE{Pille})
erhöht zusammen mit Zigarettenkonsum das Thromboserisiko
drastisch.
Im Rahmen anderer Erkrankungen, z.\,B. Tumorerkankungen, kann
das Gerinnungsystem gestört sein, wodurch es zu Thrombosen
kommen kann.
}{
\begin{itemize}
  \item Venöse Insuffizienz (Krampfadern), Rückstau
  \item Immobilisation, Bettlägrigkeit
  \item Lange Reisen ohne Bewegung: (\T{Economy class Syndrom})
  \item Mangelnde Flüssigkeitsaufnahme
  \item Rauchen, Rauchen + Pille
  \item Andere Erkrankungen
\end{itemize}
}{}{}{}

\Layout{0}{\TxSymptome}
{
Leitsymptom ist eine \EE{schmerzhafte Extremität}, zusammen
mit einer (meist einseitigen) \EE{Schwellung},
\EE{Überwärmung} und einer dunkel-bläulich-violetten
(lividen) \EE{Verfärbung}.

In der Regel ist nur eine Extremität betroffen, d.h. die
Symptome sind nur auf \E{einer Seite} zu finden.
\bigskip

% \bigskip
}{
\begin{compactitem}
  \item Schmerzen, Schwellung
  \item Dunkel-bläuliche (livide) Verfärbung. 
  \item Überwärmung
  \item Symptome i.\,d.\,R. \E{einseitig}
\end{compactitem}
}{}{}{}

\Layout{0}{\TxGefahren}
{
Es kann zu einer \E{Loslösung von Teilen des Thrombus}
kommen, dieser losgelöste Teil wird
\T{Embolus}\index{Embolus} genannt. Mit dem Blutstrom
gelangt der Embolus über die Hohlvenen in das rechte Herz.
Bis hierhin gibt es noch keine Probleme, da die Blutgefäße
an Dicke zunehmen. Vom rechten Herzen gelangt der
Embolus in den Lungenkreislauf (Lungenarterien). Hier
verzweigen sich die Geäße wieder und der Durchmesser
wird kleiner. Dies führt dazu, dass der Embolus irgendwann
zu groß ist und stecken bleibt. Dadurch wird nun ein 
Lungengefäß verstopft und ein Teil der Lunge nimmt nicht
mehr am Gasaustausch teil: Es entsteht eine
\EE{Lungenembolie}
(\prettyref{topic:lungenembolie}).
}{
\begin{itemize}
  \item Lösung des Thrombus →
  \EE{Lungenembolie} (\prettyref{topic:lungenembolie})
\end{itemize}
}{}{}{}




\begin{PictureSeriesWide}{}{Bilderserie: Thromboserisiken}
\PictureSeriesRowWide{3}
{./images/economyclass.jpg}
{Die Economyclass. Sorgsam geschlichtet verbringen Menschen
hier Stunden damit, Thrombosen zu basteln.
\SourcePicture{Gabriel}}
{./images/lovenox1.jpg}
{Throm\-bose\-prophy\-laxe. \Z{Niedermolekulares Heparin
(hier \enquote{Lovenox\texttrademark}) verhindert Thrombosen die z.\,B. durch lange
Immobilisation (Reisen, Bett\-lägrig\-keit, Gips, \ldots)
entstehen können. Die Substanz wird unter die Haut (\enquote{subkutan}) gespritzt.
\SourcePicture{Gabriel}}}
{./images/hirtler-aass/Thrombus-Embolie}
{Eine Venenthrombose kann eine Embolie auslösen. 
\SourcePicture{Hirtler}}
{empty}
{}
\end{PictureSeriesWide}


\MassnahmenMK{Spezielle
Maßnahmen}{Beinvenenthrombose}{m:beinvenenthrombose}{
\EE{Taktik}: \Speech{Symptomatische Therapie}}{}


% \MassnahmenNG{Spezielle
% Maßnahmen}{Beinvenenthrombose}{\label{m:beinvenenthrombose}}{
% \EE{Taktik}: \Speech{Symptomatische Therapie}
% 
% \begin{itemize}
%   \item Lagerung:  liegend, betroffenes Bein hochlagern
%   \item Bewegungsverbot (Emboliegefahr!)
%   \item Auf Zeichen einer Lungenembolie achten (\FullPrettyRef{topic:lungenembolie})
% \end{itemize}
% }



% \dtldefaultkey{1}
\ProcedureTest{VENGEVIX}

\citep{Renz-Polster2006,SiewertChirurgie8,ScholzNotfallmedizin2,IntroClinEmergMed4,Longmore2007,Fauci2008}

% \paragraph{Weiters}
\nocite{Eisenburger2008,
Havel2007,
AMLS3,
Erc2005DeSec5,
Erc2005Sec5,
Erc2005DeSec6,
Erc2005DeSec3,
Erc2005Sec3,
Juurlink2005,
Kalla2006,
Vlcek2008,
Wells2000,
BoehmerTaschRett6,
LPN3,
Boecker2,
Helfen2008,
LippertAnatomie7,
Patzelt2005,
Sachsenweger2003,
RedelsteinerHRNS1,
ForMed2,
StriebelAn7,
Zeiler2001,
KlinkePhysio3,
DRAn3,
ForthPharma8,
LuellmannPharma16,
ThierbachInterhospitaltransfer1,
AnCompact3,
RueckerBildatlas1}

%\putbib
% % Rea während Transport
% \nocite{Eisenburger2008,Havel2007}
% % General
% \nocite{SiewertChirurgie8}
% \nocite{AMLS3}
% \nocite{Erc2005DeSec5}
% \nocite{Erc2005Sec5}
% \nocite{Erc2005DeSec6}
% \nocite{Erc2005DeSec3}
% \nocite{Erc2005Sec3}
% \nocite{Juurlink2005}
% \nocite{Kalla2006}
% 
% \nocite{Vlcek2008}
% \nocite{Wells2000}
% \nocite{BoehmerTaschRett6}
% \nocite{Fauci2008}
% \nocite{Lissauer2007}
% 
% \nocite{LPN3,Boecker2}
% \nocite{Helfen2008}
% \nocite{LippertAnatomie7}
% \nocite{Nilsson1997}
% \nocite{Patzelt2005}
% \nocite{Sachsenweger2003}
% \nocite{RedelsteinerHRNS1}
% \nocite{ForMed2}
% \nocite{StriebelAn7}
% \nocite{Zeiler2001}
% \nocite{KlinkePhysio3}
% \nocite{DRAn3}
% \nocite{ForthPharma8}
% \nocite{LuellmannPharma16}
% \nocite{ThierbachInterhospitaltransfer1}
% \nocite{AnCompact3}
% \nocite{RueckerBildatlas1}


% \clearpage
\ChapterBibliography 

%\bibliographystyle{unsrtdin2}{\scriptsize \bibliography{Literatur-AAS}}

\ChapterEnd